\documentclass[10pt,a4paper]{amsart}
\usepackage[margin=19mm]{geometry}
\usepackage{amsmath,amssymb,mathtools}
\usepackage[scr=rsfs,cal=euler]{mathalfa}
\usepackage{xfrac}
\usepackage[unicode,hidelinks,bookmarks=false]{hyperref}
\newcommand{\ie}{i.e.\ }
\newcommand{\cf}{cf.\ }
\newcommand{\resp}{respectively\ }
\newcommand{\Subsection}{\subsection}
\providecommand{\nobreakdash}{\nobreak\hbox{-}}
\setlength{\emergencystretch}{2em}
\multlinegap0pt
\usepackage{amssymb}
\usepackage{mathtools}
\newtheorem{theorem}{Theorem}
\newtheorem{lemma}{Lemma}
\newcommand{\gammatwo}{\gamma_{2}}
\title{The 2-Domination Number and the Upper Median Degree:\\ A Proof of Graffiti.pc Conjecture 387}
\author{Jun Qing}
\date{Source: arXiv:2607.27246v1 (CC BY 4.0)}
\begin{document}
\maketitle

\noindent\emph{Source and licence.} The 2-Domination Number and the Upper Median Degree:\\ A Proof of Graffiti.pc Conjecture 387. Version arXiv:2607.27246v1 (CC BY 4.0), distributed by arXiv under the Creative Commons Attribution 4.0 International licence, http://creativecommons.org/licenses/by/4.0/, as recorded in the arXiv licence metadata for this exact version and shown on the abstract page https://arxiv.org/abs/2607.27246v1. Full licence notice: \texttt{licenses/arxiv-2607.27246-CC-BY-4.0.txt}.\par\smallskip

\noindent\textit{Editorial scope.} Counted unit: the article's theorem thm:main (source0.tex lines 87--92). The article's complete original proof of it is delivered with it (source0.tex lines 160--282, one \texttt{proof} environment of 123 lines, which the article places away from the statement). No mathematics is written, repaired or re-derived: the statement and the proof are the article's own lines, in the article's own notation, and no retained line is substituted or re-flowed.  Retained uncounted as the source-local context the proof needs: lines 99--119.  Nothing was omitted from any retained span, so the package carries no omission record.  No retained line contains a citation, so the wrapper carries no bibliography and no bibliography entry.  No retained line contains a figure environment, an image-inclusion command or drawing code, so no image is delivered and the package declares no asset dependency.  Editorial formatting only, in addition: an AMS \texttt{amsart} preamble that restates the article's own declarations and macros used by the retained lines, namely the environments \texttt{theorem}, \texttt{lemma} on separate counters, together with the standard packages those lines need.  The retained environments print in this document as Theorem 1, Lemma 1 in order of appearance, whereas the article prints the same items with the numbering of its own full document.  The complete native source of the article as distributed in the arXiv e-print for this exact version is bundled unchanged as \texttt{sources/206341/source0.tex}.
\begin{lemma}\label{lem:polynomial}
Let $B$ be a finite set, let $t\geq 0$, and choose pairwise distinct real
numbers $a_z$ for $z\in B$.  Suppose $A_1,\ldots,A_r$ are $t$-element
subsets of $B$, and define
\[
  p_i(X)=\prod_{z\in A_i}(X-a_z) \qquad (1\leq i\leq r).
\]
Assume that the indexed family $(p_i)_{i=1}^{r}$ is minimally linearly
dependent over $\mathbb{R}$.  Put
\[
  P=\bigcap_{i=1}^{r}A_i,
  \qquad s=|P|.
\]
Then:
\begin{enumerate}
  \item $2\leq r\leq t+2$;
  \item every element of $B\setminus P$ belongs to at most $r-2$ of the sets
        $A_1,\ldots,A_r$;
  \item $s\leq t+2-r$.
\end{enumerate}
\end{lemma}

\begin{theorem}\label{thm:main}
Let $G$ be a nonempty finite simple graph of order $n$.  Then
\[
  \gammatwo(G)\leq n-m(G)+1.
\]
\end{theorem}

\begin{proof}[Proof of Theorem~\ref{thm:main}]
Write $m=m(G)$.  If $m=0$, then $V(G)$ is a 2-dominating set and
\[
  \gammatwo(G)\leq n<n+1=n-m+1.
\]
Hence assume $m\geq 1$.  Let $F=\overline G$ and put
\[
  t=n-1-m,
  \qquad q=t+2=n-m+1.
\]
In particular, $0\leq t\leq n-2$ and $2\leq q\leq n$.

Define
\[
  H=\{x\in V(G):d_F(x)\leq t\},
  \qquad B=V(G)\setminus H,
\]
and write $h=|H|$ and $k=|B|$.  Since
\[
  d_F(x)=n-1-d_G(x),
\]
a vertex lies in $H$ exactly when its degree in $G$ is at least $m$.  By the
definition of the upper median, at least $\lceil n/2\rceil$ vertices have
degree at least $m$.  Thus
\begin{equation}\label{eq:majority}
  h\geq\left\lceil\frac n2\right\rceil,
  \qquad
  k\leq\left\lfloor\frac n2\right\rfloor,
  \qquad h\geq k.
\end{equation}

First suppose that $k\leq t+1$.  Since $k<q\leq n$, choose a $q$-element set
$D$ with $B\subseteq D$.  Every vertex $v\notin D$ belongs to $H$, so
\[
  d_F(v,D)\leq d_F(v)\leq t.
\]
As $v\notin D$, its neighbors in $F$ among the vertices of $D$ are exactly
its nonneighbors in $G$ among those vertices.  Hence
\[
  d_G(v,D)=|D|-d_F(v,D)\geq q-t=2.
\]
Thus $D$ is 2-dominating.

It remains to consider the case
\begin{equation}\label{eq:largeB}
  k\geq t+2.
\end{equation}
For each $x\in H$, we have
\[
  |N_F(x)\cap B|\leq d_F(x)\leq t.
\]
Using~\eqref{eq:largeB}, extend $N_F(x)\cap B$ to a $t$-element subset
$A_x\subseteq B$.  Choose pairwise distinct real numbers $a_z$ for $z\in B$,
and define
\[
  p_x(X)=\prod_{z\in A_x}(X-a_z) \qquad (x\in H),
\]
where the empty product is $1$ when $t=0$.

The vector space of real polynomials of degree at most $t$ has dimension
$t+1$.  By~\eqref{eq:majority} and~\eqref{eq:largeB},
\[
  h\geq k\geq t+2,
\]
so the indexed family $(p_x)_{x\in H}$ is linearly dependent.  Choose an
inclusion-minimal dependent subfamily, indexed by a set $C\subseteq H$, and
put $r=|C|$.  Apply Lemma~\ref{lem:polynomial} to the sets $(A_x)_{x\in C}$.
With
\[
  P=\bigcap_{x\in C}A_x,
  \qquad s=|P|,
\]
the lemma gives
\begin{equation}\label{eq:intersection}
  s\leq t+2-r,
\end{equation}
and every element of $B\setminus P$ belongs to at most $r-2$ of the sets
$A_x$ with $x\in C$.

Because $2\leq r\leq t+2$, we have $0\leq t+2-r\leq t$.  Together with
\eqref{eq:intersection} and $|B|=k\geq t+2$, this allows us to choose
$Y\subseteq B$ such that
\[
  P\subseteq Y,
  \qquad |Y|=t+2-r.
\]
Set
\[
  D=C\cup Y.
\]
Since $C\subseteq H$ and $Y\subseteq B$, the union is disjoint and
\[
  |D|=r+(t+2-r)=t+2=q.
\]

Let $v\notin D$.  If $v\in H$, then simply
\[
  d_F(v,D)\leq d_F(v)\leq t.
\]
Now suppose $v\in B$.  Since $P\subseteq Y$ and $v\notin Y$, we have
$v\notin P$.  Lemma~\ref{lem:polynomial} therefore shows that $v$ belongs to
at most $r-2$ of the sets $A_x$ with $x\in C$.  As
$N_F(x)\cap B\subseteq A_x$ for every $x\in C$, and $F$ is undirected,
\[
  d_F(v,C)
  =|\{x\in C:v\in N_F(x)\cap B\}|
  \leq r-2.
\]
Consequently,
\[
  d_F(v,D)=d_F(v,C)+d_F(v,Y)
  \leq (r-2)+|Y|=t.
\]
We have shown that every $v\notin D$ satisfies $d_F(v,D)\leq t$.  Therefore
\[
  d_G(v,D)=|D|-d_F(v,D)\geq(t+2)-t=2.
\]
Thus $D$ is a 2-dominating set of cardinality
\[
  |D|=q=n-m+1,
\]
and the theorem follows.
\end{proof}

\end{document}
